% python-torque.tex
\documentclass[11pt, a4paper]{article}

\usepackage[T1]{fontenc}
\usepackage[utf8]{inputenc}
\usepackage{tgpagella}
\usepackage{tgcursor}
\usepackage{microtype}
\usepackage[spanish, es-noshorthands]{babel}
\usepackage[margin=2.5cm]{geometry}
\usepackage{fancyhdr}
\usepackage[most]{tcolorbox}
\usepackage{listings}

\lstset{
    language=Python,
    basicstyle=\ttfamily\small,
    backgroundcolor=\color{gray!5},
    frame=single,
    rulecolor=\color{gray!40},
    keywordstyle=\color{blue}\bfseries,
    stringstyle=\color{red!70!black},
    commentstyle=\color{green!60!black},
    showstringspaces=false,
    breaklines=true,
    breakatwhitespace=true,
    columns=flexible
}

\pagestyle{fancy}
\fancyhf{}
\setlength{\headheight}{16pt}
\lhead{\textbf{UCB Sede Tarija} | Ing. Mecatrónica}
\chead{Código: Torque y Trayectoria 1R}
\rhead{Opcional W09-S02}

\definecolor{ucbblue}{RGB}{0, 51, 102}

\begin{document}

\begin{center}
    \Large\textbf{\textcolor{ucbblue}{Verificación Computacional: Perfil de Torque LSPB}}\\[0.2cm]
\end{center}

Este script conecta la trayectoria LSPB de la sesión anterior con la evaluación dinámica de torque a lo largo del tiempo, aislándolo por contribuciones físicas[cite: 14].

\begin{lstlisting}
import numpy as np
import matplotlib.pyplot as plt

# Modelo 1R Ilustrativo y Pedagogico
J = 0.25
m = 1.0
lc = 0.5
g = 9.81

# Parametros LSPB
qi = 0.0
qf = 2.0
tf = 3.0
acc = 1.0
tc = 1.0

t = np.linspace(0.0, tf, 601)

q = np.zeros_like(t)
qd = np.zeros_like(t)
qdd = np.zeros_like(t)

# Generacion del Perfil Cinetico
for k, tk in enumerate(t):
    if tk <= tc:
        q[k] = qi + 0.5 * acc * tk**2
        qd[k] = acc * tk
        qdd[k] = acc
    elif tk <= tf - tc:
        q[k] = qi + acc * tc * (tk - tc/2)
        qd[k] = acc * tc
        qdd[k] = 0.0
    else:
        q[k] = qf - 0.5 * acc * (tf - tk)**2
        qd[k] = acc * (tf - tk)
        qdd[k] = -acc

# Computo del Esfuerzo Dinamico (Torques)
tau_inertia = J * qdd
tau_velocity = np.zeros_like(t) # C=0 para 1R escalar
tau_gravity = m * g * lc * np.sin(q)
tau_total = tau_inertia + tau_velocity + tau_gravity

print("Pico absoluto |tau|:", np.max(np.abs(tau_total)))

# NOTA: Los estudiantes deberan requerir graficas de:
# 1. q(t), qdot(t), qddot(t)
# 2. tau_inertia vs t
# 3. tau_gravity vs t
# 4. tau_total vs t
\end{lstlisting}

\begin{tcolorbox}[colback=white, colframe=ucbblue, title=\textbf{Interpretación Obligatoria}]
Al graficar los datos calculados, el estudiante debe notar que:
\begin{itemize}
    \item El torque inercial (\texttt{tau\_inertia}) es un reflejo exacto del perfil discontinuo de aceleración $\ddot{q}$[cite: 14].
    \item El torque de gravedad (\texttt{tau\_gravity}) evoluciona suavemente como el seno de la posición[cite: 14].
    \item El torque total no alcanza necesariamente su pico durante la máxima aceleración, ya que la contribución gravitacional puede sumarse o restarse dependiendo de la configuración espacial del robot en ese instante[cite: 14].
\end{itemize}
\end{tcolorbox}

\end{document}
